\subsection{Berechnung der Stromverstärkung einer Transistorschaltung}
% Alt: Aufgabe 18
\Aufgabe{
    \label{Aufg18}
    Gegeben ist die folgende Schaltung mit dem Kennlinienfeld des Transistors. Die Versorgungsspannung beträgt \( U_\mathrm{V} = \mathrm{6\,V} \). Der Arbeitspunkt liegt bei der halben Betriebsspannung.\\

    \begin{figure}[H]
        \begin{minipage}[c]{\textwidth}
            \centering
            \input{Tikz/src/FigTransistorStromverstaerkung} \alttext{Die Abbildung zeigt eine elektrische Schaltung und zwei Diagramme. Die elektronische Schaltung besteht aus Transistor in der Mitte. Links vom Transistor sind zwei Widerstände, \(R_2\) oben und \(R_1\) unten, die beide mit der Basis des Transistors verbunden sind. \(R_2\) und \(R_1\) wirken dabei als Spannungsteiler. Der Widerstand \(R_2\) ist außerdem mit einer positiven Spannungsquelle \(U_\mathrm{V}\) verbunden. Der Widerstand \(R_1\) ist mit Masse (Erdung) verbunden. Am linken Rand befindet sich ein Eingangspunkt \(U_\mathrm{E}\), der ebenfalls mit Masse verbunden ist. Rechts vom Transistor ist ein Ausgangspunkt \(U_A\), der nach unten zur Masse führt. Der Emitter des Transistors ist direkt mit Masse verbunden, während der Kollektor über den Widerstand \(R_\mathrm{C}\) mit der positiven Spannungsquelle \(U_\mathrm{V}\) verbunden ist.}
        \end{minipage}
        \vspace{0.5cm} % Vertikaler Abstand zur zweiten Reihe
    \end{figure}

    \begin{figure}[H]
        \begin{minipage}[c]{0.49\textwidth}
            \centering
            \includesvg[scale=.4]{Bilder/Aufgaben/FigTransistorStromverstaerkung2} \alttext{Die zwei Diagramme können wie folgt beschrieben werden. Links ist ein Graph mit der Beschriftung auf der y-Achse „\(I_\mathrm{C}\) in \(\mathrm{mA}\)“ und auf der x-Achse „\(U_{\mathrm{CE}}\) in \(\mathrm{V}\)“. Die y-Achse reicht von \(0\) bis \(10 \ \mathrm{mA}\), die x-Achse von \(0\) bis \(6 \ \mathrm{V}\). Es sind fünf Kurven eingezeichnet, die alle bei etwa \(0 \ \mathrm{V}\) starten und schnell ansteigen, dann flacher verlaufen. Jede Kurve ist mit einem Wert für \(I_\mathrm{B}\) in Milliampere beschriftet: 5, 10, 15, 20, 25 und oben rechts „\(I_\mathrm{B} = 30 \ \mathrm{μA} \) “}
        \end{minipage}
        \hfill
        \begin{minipage}[c]{0.49\textwidth}
            \centering
            \includesvg[scale=.4]{Bilder/Aufgaben/FigTransistorStromverstaerkung3} \alttext{Rechts ist ein weiteres Diagramm mit der y-Achse beschriftet mit „\(I_\mathrm{B}\) in \(\mathrm{μA}\)“ von \(0\) bis \(20 \ \mathrm{μA}\) und die x-Achse mit „\(U_\mathrm{BE}\) in \(\mathrm{V}\)“ von \(0\) bis knapp unter \(1 \ \mathrm{V}\). Es ist eine einzelne steile Kurve eingezeichnet, die bei ca. \(0,55 \ \mathrm{V}\) stark ansteigt. zwischen \(0\) und \(0,4 \ \mathrm{V}\) ist die Kurve annähernd linear.}
        \end{minipage}
        \label{fig:FigTransistorStromverstaerkung}
    \end{figure}

    Bestimmen Sie die folgenden Werte aus dem Kennlinienfeld:

    \begin{itemize}
        \item Die Paare \( I_\mathrm{B} / U_\mathrm{BE} \) und \( I_\mathrm{C} / U_\mathrm{CE} \) im Arbeitspunkt
        \item Die Stromverstärkung \( \beta \)
    \end{itemize}
}


\Loesung{
    \begin{figure}[H]
        \centering
        \begin{subfigure}[b]{0.45\textwidth}
            \centering
            \includesvg[scale=.4]{Bilder/Aufgaben/LsgTransistorStromverstaerkung1}

        \end{subfigure}\label{fig:LsgTransistorStromverstaerkung1}
        \begin{subfigure}[b]{0.45\textwidth}
            \centering
            \includesvg[scale=.4]{Bilder/Aufgaben/LsgTransistorStromverstaerkung2}

        \end{subfigure}\label{fig:LsgTransistorStromverstaerkung2}
    \end{figure}
    Gegeben sind folgende Werte im Arbeitspunkt:
    \begin{gather*}
        U_\mathrm{CC} = \mathrm{6\,V}      \\
        U_\mathrm{CE} = \mathrm{3\,V}      \\
        I_\mathrm{C} = \mathrm{6\,mA}      \\
        U_\mathrm{BE} = \mathrm{0{,}62\,V} \\
        I_\mathrm{B} = \mathrm{20\,\mu A}
    \end{gather*}

    Berechnung der Stromverstärkung \( \beta \):
    \begin{equation*}
        \beta = \frac{I_\mathrm{C}}{I_\mathrm{B}} = \frac{\mathrm{6\,mA}}{\mathrm{20\,\mu A}} = \mathrm{300}
    \end{equation*}

    Siehe: Abschnitt \ref{sec:ElektrischesVerhaltenBipolartransistor}
}